\chapter{Conclusion}
\label{ch:conclusion}

% ─────────────────────────────────────────────────────────────────────────────
\section{Summary of Contributions}
\label{sec:conc_contributions}
% ─────────────────────────────────────────────────────────────────────────────

This project set out to design, implement, and evaluate a hybrid IRT 3PL +
BKT adaptive intelligent tutoring system for UK KS3/GCSE Geography. Five
concrete contributions have been made.

\textbf{First}, a production-deployed adaptive ITS for UK Geography was
built and made publicly accessible, filling a gap in a domain that has received
almost no attention in the adaptive learning literature \parencite{bednarz2013}.
The system handles real user accounts, persists learner state across sessions,
and serves live question content to any learner with an internet connection.

\textbf{Second}, a novel credit-factor modification to the standard BKT
learning transition equation was designed and implemented, formalising the
evidential relationship between scaffolding usage and mastery evidence. To the
author's knowledge, this specific implementation---interpolating between the
full Bayesian posterior and the prior in proportion to the credit
factor---has not been published in the existing BKT literature.

\textbf{Third}, Bloom's taxonomy was implemented as a computational selection
constraint rather than a difficulty label, enforcing that learners consolidate
recall-level knowledge before encountering understanding and application items.

\textbf{Fourth}, a five-track multi-method evaluation framework was developed
for small-sample ITS validation. This framework is reusable: any ITS project
facing insufficient sample size for inferential statistics can adopt the same
combination of algorithmic validation, engagement data, usability metrics,
misconception analysis, and robustness monitoring.

\textbf{Fifth}, a domain-agnostic open-source ITS architecture was produced.
The adaptive engine contains no Geography-specific logic; the codebase is
publicly available and provides a reusable implementation foundation for
researchers and developers building adaptive systems for declarative-knowledge
non-STEM domains.

% ─────────────────────────────────────────────────────────────────────────────
\section{Research Question Revisited}
\label{sec:conc_rq}
% ─────────────────────────────────────────────────────────────────────────────

The research question asked whether a hybrid IRT 3PL + BKT adaptive engine,
implemented in a production web application, can deliver personalised question
sequencing that produces measurable within-session learning gains for UK
KS3/GCSE Geography learners. The answer is: \textit{yes, with qualification}.
The algorithmic and technical components of the claim are fully established.
The empirical component is supported at proof-of-concept level by the $n = 4$
evaluation study, but cannot be confirmed at statistical confidence without a
powered study with a completed control condition. Overclaiming empirical
evidence is the most common failure mode in ITS evaluation research
\parencite{woolf2010}; this is an honest characterisation of the evidence.

% ─────────────────────────────────────────────────────────────────────────────
\section{Broader Significance}
\label{sec:conc_significance}
% ─────────────────────────────────────────────────────────────────────────────

\textbf{ITS for non-STEM domains.} This project demonstrates that the IRT+BKT
hybrid architecture transfers to a declarative-knowledge domain with a
multi-dimensional prerequisite structure, and that ZPD targeting, mastery
gating, and spaced relearning are applicable to Geography content. This extends
the empirical and conceptual scope of ITS research into territory that has been
largely neglected.

\textbf{Open-source adaptive learning infrastructure.} Commercial adaptive
platforms are black boxes: algorithms proprietary, calibration data unavailable,
architectures inextensible. This project provides a fully transparent, auditable
alternative where every adaptive decision can be traced to specific parameter
values in the open-source codebase.

\textbf{Domain-agnostic architecture as a research platform.} The most
significant future direction, Cross-Domain Knowledge Transfer, is
architecturally feasible without algorithmic changes. This positions the project
not as a finished product but as a research platform for investigating transfer
learning in adaptive educational systems.

% ─────────────────────────────────────────────────────────────────────────────
\section{Personal Reflection}
\label{sec:conc_reflection}
% ─────────────────────────────────────────────────────────────────────────────

The scope reduction at Sprint~5 was the most professionally significant
decision of the project. The original ambition was larger; the final system
is more rigorous. That trade-off---depth over breadth, defensible claims over
impressive-sounding ones---is the central lesson of the project. The experience
of implementing a psychometric algorithm in a production web application has
provided a level of engineering judgement that no coursework exercise could
replicate: the Review Material bug, the EAP grid resolution issue, the K-means
supersession are not failures but the data points from which engineering
judgement is built.

% ─────────────────────────────────────────────────────────────────────────────
\section{Final Statement}
\label{sec:conc_final}
% ─────────────────────────────────────────────────────────────────────────────

Bloom's 2-sigma problem has not been solved. Human tutoring remains superior
to any automated system across the full range of educational contexts and
learner profiles. But the problem is tractable: the literature shows that
well-designed adaptive systems can close the gap substantially, and that the
key design variables are interaction granularity, principled algorithm selection,
and transparent feedback, not algorithmic complexity or scale.

GeoMentor is a small contribution to a large problem. It demonstrates that
these design principles can be implemented at production quality by a single
developer, in a non-STEM domain, using open-source tools, within an academic
project timeline. If it serves as a foundation for a powered evaluation study,
a secondary school pilot, or a researcher building an adaptive system for
History or Computer Science, it will have served its purpose.
